\chapter{Programs — Pass}
%%% generated by the blueprint pipeline from TCSlib.Programs.Light.Lib.Pass
\section{Overview}

Declarations of the Lean module \texttt{TCSlib.Programs.Light.Lib.Pass}.

\section{Declarations}

\begin{definition}[Indexed store pass]
\lean{Light.pass}
\label{Light.pass}
\leanok
\uses{Light.Expr, Light.Stmt, Light.Stmt.for, Light.v}
For natural numbers $c$ and a destination expression $d$, a length expression $\ell$,
and a value expression $e$, the pass is the statement that loops over $c$ up to $\ell$
and stores $e$ at the address $d+c$ on each iteration.
\end{definition}

\begin{theorem}[A Written Segment Matches Its List]
\lean{Light.seg_wrote}
\label{Light.seg_wrote}
\leanok
\uses{Light.Seg, Light.wrote, Light.wrote_done}
\difficulty{3}
Let $\mu,f:\mathbb{N}\to\mathbb{Z}$, let $d,n\in\mathbb{N}$, and let $L$ be a finite
list of integers with length $n$. Suppose that $L_i=f(i)$ for every valid index $i$.
After writing the first $n$ values of $f$ into $\mu$ starting at $d$, the resulting
sequence represents $L$ starting at $d$.
\end{theorem}

\begin{theorem}[A Mapped List in a Written Segment]
\lean{Light.seg_wrote_map}
\label{Light.seg_wrote_map}
\leanok
\uses{Light.Seg, Light.seg_wrote, Light.wrote}
\difficulty{3}
Let $\mu:\mathbb{N}\to\mathbb{Z}$, let $d\in\mathbb{N}$, let
$g:\mathbb{Z}\to\mathbb{Z}$, and let $L$ be a finite list of integers. Write the values
$g(L_i)$ into $\mu$ starting at $d$, using the offset $i$ for each entry. The resulting
sequence represents the list obtained by applying $g$ to each entry of $L$, starting at
$d$.
\end{theorem}

\begin{theorem}[Correctness of an indexed store pass]
\lean{Light.Ends.pass}
\label{Light.Ends.pass}
\leanok
\uses{Light.Ends, Light.Ends.block, Light.Ends.forMem, Light.Expr, Light.Expr.Safe, Light.Expr.cost, Light.Expr.val, Light.Limits.Addr, Light.Op.eval, Light.State, Light.Stmt.after, Light.Stmt.blockCost, Light.pass, Light.toNat_natCast_add_natCast, Light.v, Light.wrote, Light.wrote_succ, Light.wrote_zero}
\difficulty{5}
Fix execution limits and a program, natural numbers $c,x,y,n,T$, an expression $e$,
integer-valued functions $f$ and $loc$, an initial memory $\mu$, a destination $dst$,
and a state predicate $Q$. For every $j<n$, suppose $e$ is safe in the state whose
locations are $loc$ updated to give variable $c$ value $j$ and whose memory is obtained
by writing the first $j$ values of $f$ starting at $dst$, and that $e$ evaluates there
to $f(j)$. Suppose $Q$ holds in the state with variable $c$ set to $n$ and memory
containing the first $n$ values of $f$ at $dst$; the space limit is at most the word
limit; $dst+n$ is at most the space limit; $loc(x)=n$; $loc(y)=dst$; $x\ne c$ and $y\ne
c$; and $n(\operatorname{cost}(e)+12)+6\le T$. Then the indexed store pass using counter
$c$, length expression $x$, destination expression $y$, and value expression $e$
executes from locations $loc$ and memory $\mu$ within cost $T$, ending in a state
satisfying $Q$.
\end{theorem}
